\documentclass[10pt,a4paper]{amsart}
\usepackage{amsmath,amssymb,amsthm,graphicx}
\usepackage{verbatim}
\newenvironment{reference}{\comment}{\endcomment}
%\newenvironment{reference}{}{}
\newenvironment{slogan}{\comment}{\endcomment}
\newenvironment{history}{\comment}{\endcomment}

% For commutative diagrams we use Xy-pic
\usepackage[all]{xy}

% We use 2cell for 2-commutative diagrams.
\xyoption{2cell}
% Theorem environments.
%
\theoremstyle{plain}
\newtheorem{theorem}[subsection]{Theorem}
\newtheorem{proposition}[subsection]{Proposition}
\newtheorem{lemma}[subsection]{Lemma}

\theoremstyle{definition}
\newtheorem{definition}[subsection]{Definition}
\newtheorem{example}[subsection]{Example}
\newtheorem{exercise}[subsection]{Exercise}
\newtheorem{situation}[subsection]{Situation}

\theoremstyle{remark}
\newtheorem{remark}[subsection]{Remark}
\newtheorem{remarks}[subsection]{Remarks}

\numberwithin{equation}{subsection}

% Macros
%
\def\lim{\mathop{\mathrm{lim}}\nolimits}
\def\colim{\mathop{\mathrm{colim}}\nolimits}
\def\Spec{\mathop{\mathrm{Spec}}}
\def\Hom{\mathop{\mathrm{Hom}}\nolimits}
\def\Ext{\mathop{\mathrm{Ext}}\nolimits}
\def\SheafHom{\mathop{\mathcal{H}\!\mathit{om}}\nolimits}
\def\SheafExt{\mathop{\mathcal{E}\!\mathit{xt}}\nolimits}
\def\Sch{\mathit{Sch}}
\def\Mor{\mathop{\mathrm{Mor}}\nolimits}
\def\Ob{\mathop{\mathrm{Ob}}\nolimits}
\def\Sh{\mathop{\mathit{Sh}}\nolimits}
\def\NL{\mathop{N\!L}\nolimits}
\def\CH{\mathop{\mathrm{CH}}\nolimits}
\def\proetale{{pro\text{-}\acute{e}tale}}
\def\etale{{\acute{e}tale}}
\def\QCoh{\mathit{QCoh}}
\def\Ker{\mathop{\mathrm{Ker}}}
\def\Im{\mathop{\mathrm{Im}}}
\def\Coker{\mathop{\mathrm{Coker}}}
\def\Coim{\mathop{\mathrm{Coim}}}

% Boxtimes
%
\DeclareMathSymbol{\boxtimes}{\mathbin}{AMSa}{"02}

%
% Macros for moduli stacks/spaces
%
\def\QCohstack{\mathcal{QC}\!\mathit{oh}}
\def\Cohstack{\mathcal{C}\!\mathit{oh}}
\def\Spacesstack{\mathcal{S}\!\mathit{paces}}
\def\Quotfunctor{\mathrm{Quot}}
\def\Hilbfunctor{\mathrm{Hilb}}
\def\Curvesstack{\mathcal{C}\!\mathit{urves}}
\def\Polarizedstack{\mathcal{P}\!\mathit{olarized}}
\def\Complexesstack{\mathcal{C}\!\mathit{omplexes}}
% \Pic is the operator that assigns to X its picard group, usage \Pic(X)
% \Picardstack_{X/B} denotes the Picard stack of X over B
% \Picardfunctor_{X/B} denotes the Picard functor of X over B
\def\Pic{\mathop{\mathrm{Pic}}\nolimits}
\def\Picardstack{\mathcal{P}\!\mathit{ic}}
\def\Picardfunctor{\mathrm{Pic}}
\def\Deformationcategory{\mathcal{D}\!\mathit{ef}}
\usepackage{cleveref}
\crefname{originallemma}{original lemma}{original lemmas}
\Crefname{originallemma}{Original lemma}{Original lemmas}
\makeatletter
\newlabel{etale-lemma-finite-etale-etale-local}{{04HN}{source}{}{sourcepointer-1}{}}
\newlabel{etale-lemma-finite-etale-etale-local@cref}{{[originallemma][0][]04HN}{[1][0][]source}{}{}{}}
\makeatother
\crefname{originallemma}{original lemma}{original lemmas}
\Crefname{originallemma}{Original lemma}{Original lemmas}
\makeatletter
\newlabel{descent-lemma-descent-data-sheaves}{{02W5}{source}{}{sourcepointer-2}{}}
\newlabel{descent-lemma-descent-data-sheaves@cref}{{[originallemma][0][]02W5}{[1][0][]source}{}{}{}}
\makeatother
\crefname{originallemma}{original lemma}{original lemmas}
\Crefname{originallemma}{Original lemma}{Original lemmas}
\makeatletter
\newlabel{descent-lemma-affine}{{0245}{source}{}{sourcepointer-3}{}}
\newlabel{descent-lemma-affine@cref}{{[originallemma][0][]0245}{[1][0][]source}{}{}{}}
\makeatother
\crefname{originallemma}{original lemma}{original lemmas}
\Crefname{originallemma}{Original lemma}{Original lemmas}
\makeatletter
\newlabel{descent-lemma-descending-property-finite}{{02LA}{source}{}{sourcepointer-4}{}}
\newlabel{descent-lemma-descending-property-finite@cref}{{[originallemma][0][]02LA}{[1][0][]source}{}{}{}}
\makeatother
\crefname{originallemma}{original lemma}{original lemmas}
\Crefname{originallemma}{Original lemma}{Original lemmas}
\makeatletter
\newlabel{descent-lemma-descending-property-etale}{{02VN}{source}{}{sourcepointer-5}{}}
\newlabel{descent-lemma-descending-property-etale@cref}{{[originallemma][0][]02VN}{[1][0][]source}{}{}{}}
\makeatother
\title{Finite locally constant sheaves are exactly the sheaves represented by finite etale schemes}
\author{The Stacks Project Authors}
\date{stacks-project:etale-cohomology.tex:local-snapshot-20261011; GNU FDL 1.2 or later}
\begin{document}
\maketitle
\noindent\textit{Editorial note.} All delivered mathematical text is the author's own, copied from the pinned Stacks Project chapter file etale-cohomology.tex; no statement or proof step was rewritten, reordered or removed.\par\medskip
\noindent\textit{Editorial note.} The proof is the author's complete argument for both implications; the parenthetical about the construction of the descent datum refers to a routine transcription of an argument in the cited descent lemma, and every subsequent step of the descent is stated with its input.\par\medskip
\noindent\textit{Editorial note.} Every labelled reference occurring inside the counted spans is carried in assembly.external\_references as a non-load-bearing navigation pointer (same-chapter ones at their definition site, cross-chapter ones at the invocation inside the counted span), so no load-bearing mathematics is replaced by a stub. No bibliography citation occurs inside the counted spans.\par\medskip
\noindent\textit{Editorial note.} The checked chapter list for same-name collisions is the set of stacks chapters staged in this repository: chow, cohomology, divisors, etale-cohomology, flat, more-algebra, more-morphisms, morphisms, spaces-more-morphisms, varieties together with this card's chapter etale-cohomology.\par\medskip
\noindent\textit{Editorial note.} The preamble is the chapter's own native preamble excerpt already used for this source: the Xy-pic setup (native preamble lines 14-18) and the native theorem declarations with the mathematical macros (native preamble lines 154-217).\par\medskip
\noindent\textit{Allowed standard prerequisites (editorial).}\begin{itemize}
\item Local structure of a finite etale morphism. Etale Morphisms, Lemma etale-lemma-finite-etale-etale-local: a finite etale morphism is locally a disjoint union of isomorphisms. Shows that a represented sheaf is finite locally constant.
\item Effectivity of descent for schemes. Descent, Lemma descent-lemma-affine together with Lemmas descent-lemma-descending-property-finite and descent-lemma-descending-property-etale. Turns the constructed descent datum into a morphism of schemes and transfers finiteness and etaleness.
\item The descent datum attached to a sheaf. Descent, Lemma descent-lemma-descent-data-sheaves, whose construction is copied. Produces the descent datum from the local representing objects.
\end{itemize}\medskip
\ifdefined\hypertarget\hypertarget{sourcepointer-1}{}\fi
\noindent\textit{Source reference.} Original lemma 04HN: sources/stacks-etale-cohomology/pinned-original-source.tex, source lines 9986--9986. Navigation pointer invoked at proof line 9986; the counted span contains the pointer only and the cited result is not reproduced inside it, so the pointer is registered as a non-load-bearing navigation stub in the established form for Stacks material. No bibliography citation is involved.\par\medskip
\ifdefined\hypertarget\hypertarget{sourcepointer-2}{}\fi
\noindent\textit{Source reference.} Original lemma 02W5: sources/stacks-etale-cohomology/pinned-original-source.tex, source lines 10011--10011. Navigation pointer invoked at proof line 10011; the counted span contains the pointer only and the cited result is not reproduced inside it, so the pointer is registered as a non-load-bearing navigation stub in the established form for Stacks material. No bibliography citation is involved.\par\medskip
\ifdefined\hypertarget\hypertarget{sourcepointer-3}{}\fi
\noindent\textit{Source reference.} Original lemma 0245: sources/stacks-etale-cohomology/pinned-original-source.tex, source lines 10014--10014. Navigation pointer invoked at proof line 10014; the counted span contains the pointer only and the cited result is not reproduced inside it, so the pointer is registered as a non-load-bearing navigation stub in the established form for Stacks material. No bibliography citation is involved.\par\medskip
\ifdefined\hypertarget\hypertarget{sourcepointer-4}{}\fi
\noindent\textit{Source reference.} Original lemma 02LA: sources/stacks-etale-cohomology/pinned-original-source.tex, source lines 10016--10016. Navigation pointer invoked at proof line 10016; the counted span contains the pointer only and the cited result is not reproduced inside it, so the pointer is registered as a non-load-bearing navigation stub in the established form for Stacks material. No bibliography citation is involved.\par\medskip
\ifdefined\hypertarget\hypertarget{sourcepointer-5}{}\fi
\noindent\textit{Source reference.} Original lemma 02VN: sources/stacks-etale-cohomology/pinned-original-source.tex, source lines 10017--10017. Navigation pointer invoked at proof line 10017; the counted span contains the pointer only and the cited result is not reproduced inside it, so the pointer is registered as a non-load-bearing navigation stub in the established form for Stacks material. No bibliography citation is involved.\par\medskip
\begin{lemma}
\label{lemma-characterize-finite-locally-constant}
Let $X$ be a scheme and $\mathcal{F}$ a sheaf of sets on $X_\etale$.
Then the following are equivalent
\begin{enumerate}
\item $\mathcal{F}$ is finite locally constant, and
\item $\mathcal{F} = h_U$ for some finite \'etale morphism $U \to X$.
\end{enumerate}
\end{lemma}

\begin{proof}
A finite \'etale morphism is locally isomorphic to a disjoint union
of isomorphisms, see
\'Etale Morphisms, Lemma \ref{etale-lemma-finite-etale-etale-local}.
Thus (2) implies (1). Conversely, if $\mathcal{F}$ is finite locally
constant, then there exists an \'etale covering $\{X_i \to X\}$ such that
$\mathcal{F}|_{X_i}$ is representable by $U_i \to X_i$ finite \'etale.
Arguing exactly as in the proof of
Descent, Lemma \ref{descent-lemma-descent-data-sheaves}
we obtain a descent datum for schemes $(U_i, \varphi_{ij})$ relative to
$\{X_i \to X\}$ (details omitted). This descent datum is effective for
example by Descent, Lemma \ref{descent-lemma-affine}
and the resulting morphism of schemes $U \to X$ is finite \'etale
by Descent, Lemmas \ref{descent-lemma-descending-property-finite} and
\ref{descent-lemma-descending-property-etale}.
\end{proof}

\end{document}
